\documentclass[11pt]{article}

\usepackage[margin=1in]{geometry}
\usepackage{amsmath,amssymb,amsthm}
\usepackage{enumitem}
\usepackage{hyperref}
\usepackage{cleveref}
\usepackage{natbib}

\hypersetup{
  colorlinks=true,
  linkcolor=blue,
  citecolor=blue,
  urlcolor=blue
}

\theoremstyle{plain}
\newtheorem{theorem}{Theorem}
\newtheorem{lemma}[theorem]{Lemma}
\newtheorem{corollary}[theorem]{Corollary}

\theoremstyle{definition}
\newtheorem{definition}[theorem]{Definition}

\theoremstyle{remark}
\newtheorem{remark}[theorem]{Remark}

\newcommand{\vars}{\operatorname{vars}}
\newcommand{\size}{\operatorname{size}}
\newcommand{\lw}{\operatorname{lw}}
\newcommand{\OBDD}{\mathrm{OBDD}}
\newcommand{\decSDNNF}{\mathrm{dec}\mbox{-}\mathrm{SDNNF}}

\title{A Common-Order Quasipolynomial Simulation from dec-SDNNF to OBDD}
\author{}
\date{\today}

\begin{document}
\maketitle

\section{Purpose}

This note records a small variant of the Bollig--Buttkus simulation of structured
decision-DNNFs by OBDDs.  The usual simulation orients each vtree split
using counts that depend on the input circuit.  That is enough for simulating one
circuit, but it can produce different output variable orders for two circuits
respecting the same vtree.

The variant below fixes the orientation from the vtree alone.  Consequently,
several circuits respecting the same finite vtree can be compiled into
OBDDs using one common variable order, while preserving a quasipolynomial size
bound.

\section{Setup}

Let $T$ be a finite vtree over a finite nonempty variable set $X$.  For a node
$t$ of $T$, write $T_t$ for the subtree rooted at $t$, and define
\[
  w(t) := |\text{leaves}(T_t)|.
\]
At each internal node $t$ of $T$ with children $t_0,t_1$, call the child with
smaller weight the light child, breaking ties by a fixed convention depending
only on $T$.  The other child is the heavy child.

This orientation induces a total order $\pi_X$ on $X$: recursively list the
variables of the light child first and the variables of the heavy child second.

\begin{remark}[Normalization]
We assume the input circuit has been cleaned so that every gate reachable from
the output computes a function depending on at least one variable, except for the
constant sinks, and every AND gate has two nonconstant inputs.  Constant-only
subcircuits can be evaluated and contracted in linear time.  This keeps every
AND split nonempty and does not affect determinism, structuredness, or
asymptotic size.
\end{remark}

\section{The Construction}

Let $D$ be a normalized $\decSDNNF$ respecting $T$.  Literals, if present, are
treated as decision nodes: a leaf $x$ tests $x$ and sends the $1$-edge to the
$1$-sink and the $0$-edge to the $0$-sink; a leaf $\neg x$ does the reverse.

We first classify edges of $D$.  Let $g$ be an AND gate.  Its two inputs are
separated by some split of $T$; one input uses variables from the light side of
that split, and the other uses variables from the heavy side.  The edge from
$g$ to the light-side input is a \emph{light edge}; the edge from $g$ to the
heavy-side input is a \emph{heavy edge}.  All edges out of decision nodes are
\emph{neutral edges}.

The rest of this construction is identical to that of Beame; only the convention
of "light" and "heavy" has been altered.

For a directed path $P$ in $D$, let $S(P)$ be the set of light edges on $P$.
This can be thought of as the "stack" of remaining AND-nodes whose heavy edges 
are left to traverse.
For a node $u$ of $D$, define
\[
  S(u) := \{S(P) \mid P \text{ is a directed path from the root of } D
          \text{ to } u\}.
\]
The output graph $F$ has one node $(u,s)$ for every node $u$ of $D$ and every
$s\in S(u)$.  Its root is $(r,\emptyset)$, where $r$ is the root of $D$.

The labels of nodes in $F$ are as follows.
\begin{enumerate}[label=(\roman*),noitemsep]
  \item If $u$ is a decision node labeled by $x$, then $(u,s)$ is a decision
  node labeled by $x$.
  \item If $u$ is an AND gate, then $(u,s)$ is an unlabeled no-op node.
  \item If $u$ is a $0$-sink, then $(u,s)$ is a $0$-sink.
  \item If $u$ is a $1$-sink and $s=\emptyset$, then $(u,s)$ is a $1$-sink.
  If $u$ is a $1$-sink and $s\neq\emptyset$, then $(u,s)$ is an unlabeled
  no-op node.
\end{enumerate}

The edges of $F$ are defined by three rules.
\begin{enumerate}[label=(\arabic*),noitemsep]
  \item For each light edge $e=(u,v)$ of $D$ and each $s\in S(u)$, add
  \[
    ((u,s),(v,s\cup\{e\})).
  \]
  This records that we have entered the first half of an AND gate and must later
  continue into its heavy sibling if this light subcomputation accepts.

  This is akin to a "stack push" operation

  \item For each neutral edge $(u,v)$ of $D$ and each $s\in S(u)$, add
  \[
    ((u,s),(v,s)).
  \]
  Neutral motion does not create or discharge any pending ANDs.

  \item For each heavy edge $(u,v_h)$ of $D$, let $e=(u,v_\ell)$ be its sibling
  light edge.  For each $s\in S(u)$ and each $1$-sink $w$ reachable in
  $D_{v_\ell}$, add
  \[
    ((w,s\cup\{e\}),(v_h,s)).
  \]
  This is akin to the "stack pop". After the light input has accepted, the computation
  pops the obligation $e$ and continues into the heavy input.
\end{enumerate}

Unlabeled no-op nodes can be contracted after the construction.  They are kept
in the description because they make the bookkeeping transparent.

After contracting the unlabeled no-op nodes, the only internal nodes left are
decision nodes.  Thus the result is an FBDD. However, we can show that all paths
respect the same order $\pi_X$, hence the result is an $OBDD_{\pi_X}$

\section{Why the Order is Common}

\begin{lemma}[Order lemma]
Let $D$ be a normalized $\decSDNNF$ respecting $T$.  The OBDD produced by the
construction above respects the order $\pi_X$.
\end{lemma}

\begin{proof}
The only point to check is an AND gate $g$ decomposed at a vtree node $t$ with
children $t_0,t_1$.  By construction, all variables in the light side of $t$ are
placed before all variables in the heavy side of $t$ in $\pi_X$.  The simulation
therefore evaluates the earlier block before the later block.

Inside each block, the same argument applies recursively to lower vtree nodes.
Decision gates do not change the vtree region in which the computation is
located; they only move within the current region.  Hence along every accepting
or rejecting computation path, variables are queried in increasing $\pi_X$ order,
with variables absent from $D$ simply skipped.
\end{proof}

\begin{corollary}[Same order for all circuits respecting the same tree]
Let $D_1,\ldots,D_k$ be normalized $\decSDNNF$s respecting the
same finite vtree $T$ over variables $X$.
Then the construction compiles every $D_i$ into an $\OBDD$ respecting the same
variable order $\pi_X$.
\end{corollary}

\section{Size Bound}

Let $N=\size(D)$ and let $M$ be the number of AND gates of $D$.  Once the
light/heavy orientation is fixed, call an edge from an AND gate to its light
input a light edge.  Here $L$ means the maximum number of light edges on any
directed source-to-sink path in the original circuit graph, where a path through
an AND gate follows one of its two outgoing edges.

The construction bounds the size of the output by
\[
  O(NM^L),
\]

with the reasoning that each node is labeled $(v, s)$ for some node $v$ from $D$ and
some stack $s$. With $M$ possible choices for an AND-node (hence for a light edge), and since
s is a set of at most $L$ light edges, we obtain the above.

It remains only to bound $L$ for the vtree-based ordering.

\begin{lemma}[Light edges halve vtree weight]
For every normalized $\decSDNNF$ $D$ respecting $T$,
every directed source-to-sink path of $D$ contains at most
$\lceil \log_2 |X| \rceil + 1$ light edges.
\end{lemma}

\begin{proof}
Consider the light edges on one path, and let their corresponding AND gates be
$g_1,\ldots,g_L$.  Let $t_i=t(g_i)$ be the vtree node used by $g_i$.

This is a path in the original circuit DAG, not the later stacked OBDD.  Hence,
after the path takes the light input of $g_i$, every subsequent gate on the path
is still inside that light-input subcircuit.  Therefore, if $i<L$, then
$t_{i+1}$ lies strictly inside the light child subtree of $t_i$.  Since that
child was chosen to have minimum weight in $T$,
\[
  w(t_{i+1}) \leq \frac{w(t_i)}{2}.
\]
By induction,
\[
  w(t_L) \leq \frac{w(t_1)}{2^{L-1}} \leq \frac{|X|}{2^{L-1}}.
\]
The final AND gate $g_L$ has two nonconstant inputs after normalization, so its
decomposition node contains at least one variable of $X$.  Hence $w(t_L)\geq 1$.
Combining the inequalities gives
\[
  1 \leq \frac{|X|}{2^{L-1}},
\]
so $L \leq \log_2 |X| + 1$.
\end{proof}

\begin{theorem}[Common-order simulation]
Let $D_1,\ldots,D_k$ be normalized $\decSDNNF$ circuits respecting the same finite
vtree $T$ over variables $X$.
For each $i$, the construction above outputs an $\OBDD$ for $D_i$ respecting the
same order $\pi_X$, of size
\[
  O\!\left(N_iM_i^{\lceil \log_2 |X| \rceil + 1}\right),
\]
where $N_i=\size(D_i)$ and $M_i$ is the number of AND gates of $D_i$.
\end{theorem}

\begin{proof}
The common-order claim is the corollary above.  The size claim is the standard
copying bound $O(N_iM_i^L)$ together with the light-edge bound
$L\leq \lceil \log_2 |X| \rceil+1$.  Because the input has only decision gates
and AND gates, the output has only decision nodes after contracting no-op nodes,
so it is an OBDD.
\end{proof}

\begin{corollary}[Two size-$N$ circuits]
If $D_1$ and $D_2$ both have size at most $N$, then
Since $M_i, X\leq N$, both circuits compile to $\OBDD$s over one common order of
size
\[
  O(NN^{\log(N) + 1}) = O(N2^{\log(N)+1})
\].

\end{corollary}

\section{Takeaway}

The original circuit-dependent orientation is useful because it can give a
slightly sharper bound for one circuit.  The vtree leaf-count orientation trades
that local sharpness for a global benefit: the output order depends only on the
shared finite vtree.  For any fixed finite family of dec-SDNNF circuits
respecting that vtree, this gives one common OBDD order and keeps the simulation
quasipolynomial in the total input size.

\end{document}
